\documentclass[12pt,a4paper]{article}

%% --- Packages ---
\usepackage[utf8]{inputenc}
\usepackage[T1]{fontenc}
\usepackage{amsmath,amssymb,amsthm}
\usepackage{mathtools}
\usepackage{slashed}          % Feynman slash
\usepackage{physics}
\usepackage{graphicx}
\usepackage{booktabs}
\usepackage{hyperref}
\usepackage{xcolor}
\usepackage[margin=2.5cm]{geometry}
\usepackage{cite}
\usepackage{bbold}

\hypersetup{
  colorlinks=true,
  linkcolor=blue!70!black,
  citecolor=green!60!black,
  urlcolor=blue!70!black
}

%% --- Custom commands ---
\newcommand{\Lc}{\mathcal{L}}
\newcommand{\Dc}{\mathcal{D}}
\newcommand{\Zc}{Z[J,\eta,\bar\eta]}
\newcommand{\mphi}{m_\phi}
\newcommand{\mchi}{m_\chi}
\newcommand{\alphass}{\alpha_s}
\newcommand{\alphapp}{\alpha_p}
\newcommand{\thetarel}{\theta_{\rm relic}}
\newcommand{\sigmaT}{\sigma_T}
\newcommand{\OmegaDM}{\Omega_{\rm DM} h^2}
\newcommand{\rholambda}{\rho_\Lambda}

%%=================================================================
\begin{document}

\title{\textbf{The Dark Unification}\\[0.5em]
\large Self-Interacting Majorana Dark Matter, Path-Integral Foundations,\\
and Dark Energy as the T-Breaking Component}

\author{Omer P.\\[0.3em]
\small Independent Researcher\\
\small \href{https://github.com/0mp3s}{github.com/0mp3s}}

\date{March 2026 \quad\quad DOI: (pending Zenodo upload)}

\maketitle

\begin{abstract}
We present a unified theoretical framework for the dark sector built from a
single Lagrangian with three free parameters $(m_\chi, m_\phi, \alpha)$.
\textbf{Chapter~1} establishes the phenomenological model: a Majorana fermion
dark matter candidate $\chi$ interacting via a real scalar mediator $\phi$.
Velocity-dependent self-interaction cross sections are computed exactly via the
Variable Phase Method (VPM), and an MCMC scan over astrophysical data identifies
122 viable benchmark points consistent with relic density, Fornax globular cluster
survival, SPARC rotation curves, the Radial Acceleration Relation, and direct/indirect
detection constraints.
\textbf{Chapter~2} derives all Chapter~1 results from the generating functional
$Z[J]$ of the path integral, proving that VPM is \emph{exactly} the fluctuation
determinant of the scattering field integral (no approximation). The Coleman-Weinberg
potential, finite-$T$ Matsubara formalism, and dark QCD misalignment together predict
$H_0 = 67.4$~km/s/Mpc from $V_{eff}(\sigma_0) = \rho_\Lambda$ via the Friedmann
equation.
\textbf{Chapter~3} proposes a Dark Electromagnetic Duality: the CP phase
$\theta \equiv \sigma/f$ is promoted to a dynamical axion field, with
$y_s = y\cos(\sigma/f)$ (T-even, SIDM clustering) playing the role of $\vec{E}$
and $y_p = y\sin(\sigma/f)$ (T-odd, vacuum energy) playing the role of $\vec{B}$.
Dark energy is not a separate phenomenon; it is the T-violating projection of the
same Yukawa interaction. The universal relic angle
$\thetarel = \arctan(1/\sqrt{8}) = 19.47^\circ$ is fixed simultaneously by
SIDM and relic density constraints, with no free parameters.
\end{abstract}

\tableofcontents
\newpage

%%=================================================================
\section*{Notation and Conventions}

\begin{itemize}
  \item Metric signature: $(+,-,-,-)$
  \item Natural units: $\hbar = c = k_B = 1$ unless stated
  \item $\alpha = y^2/(4\pi)$, $\alphass = y_s^2/(4\pi)$, $\alphapp = y_p^2/(4\pi)$
  \item $\lambda = \alpha m_\chi / m_\phi$ (dimensionless coupling)
  \item $\kappa = \mu v_{\rm rel}/m_\phi$ (dimensionless momentum), $\mu = m_\chi/2$
  \item Lagrangian sign convention: $\Lc = T - V$, so $V(r) < 0$ for attraction
\end{itemize}

\newpage
%%=================================================================
%%=================================================================
\part{Secluded Majorana SIDM: Phenomenology}
\label{part:sidm}
%%=================================================================
%%=================================================================

%%=================================================================
%% CHAPTER 1: Secluded Majorana SIDM — Phenomenology
%% Source: github.com/0mp3s/Secluded-Majorana-SIDM (DOI: 10.5281/zenodo.19225823)
%%=================================================================

\section{Introduction}
\label{ch1:intro}

The small-scale structure anomalies of $\Lambda$CDM — the core-cusp problem,
missing satellites, too-big-to-fail — motivate self-interacting dark matter (SIDM),
where DM particles scatter elastically via a light mediator. We consider a
\emph{secluded} dark sector: a Majorana fermion $\chi$ coupled to a real scalar
$\phi$ with no tree-level coupling to the SM.

\section{The Model}
\label{ch1:model}

The interaction Lagrangian is:
\begin{equation}
  \Lc_{\rm int} = -\frac{y}{2}\bar{\chi}\chi\phi
\end{equation}

In the non-relativistic limit, scalar exchange generates an attractive Yukawa potential:
\begin{equation}
  V(r) = -\frac{\alpha}{r}e^{-\mphi r}, \qquad \alpha = \frac{y^2}{4\pi}
\end{equation}

The model has three free parameters: $(m_\chi,\, m_\phi,\, \alpha)$.

\subsection{Scalar Potential and Vacuum Stability}

The full scalar potential:
\begin{equation}
  V(\phi) = \frac{1}{2}m_\phi^2\phi^2 + \frac{\mu_3}{3!}\phi^3 + \frac{\lambda_\phi}{4!}\phi^4
\end{equation}

Vacuum stability (no secondary minimum deeper than $\phi=0$) requires:
\begin{equation}
  \lambda_\phi > \frac{3\mu_3^2}{8m_\phi^2}
\end{equation}

For the cannibal lower bound $\mu_3/m_\phi = 1.7$: $\lambda_\phi \gtrsim 1.08$.
Numerically: $\lambda_\phi \gtrsim 0.96$ (from $V(\phi_{\rm min}) > 0$ scan).

\section{Self-Interaction Cross Section: Variable Phase Method}
\label{ch1:vpm}

The transfer cross section is computed via the VPM ODE for phase shifts $\delta_\ell$:
\begin{equation}
  \frac{d\delta_\ell}{dx} = +\frac{\lambda\,e^{-x}}{\kappa\,x}
  \left[\hat{j}_\ell(\kappa x)\cos\delta_\ell - \hat{n}_\ell(\kappa x)\sin\delta_\ell\right]^2
\end{equation}

where $x = m_\phi r$, $\kappa = \mu v_{\rm rel}/m_\phi$, $\lambda = \alpha m_\chi/m_\phi$.

For identical Majorana fermions, quantum statistics requires:
\begin{equation}
  \sigmaT = \frac{2\pi}{k^2}\left[\sum_{l\,\rm even}(2l+1)\sin^2\delta_l
             + 3\sum_{l\,\rm odd}(2l+1)\sin^2\delta_l\right]
\end{equation}

\textbf{Key result (proved in Chapter~2, Appendix~\ref{app:vpm_proof}):} VPM is
the \emph{exact} fluctuation determinant of the path integral — no approximation.

\section{Relic Density}
\label{ch1:relic}

Freeze-out via $\chi\chi \to \phi\phi$ (t/u-channel scalar mediator).
For Majorana fermions, this is s-wave:
\begin{equation}
  \langle\sigma v\rangle_0 = \frac{\pi\alpha^2}{4m_\chi^2}
\end{equation}

Sommerfeld enhancement at freeze-out: $S_0 < 1.026$ for all 17 relic-viable BPs
(correction $< 2.5\%$). CMB constraint: inapplicable (secluded sector, $f_{\rm eff}=0$).

\section{MCMC Posterior and Viable Island}
\label{ch1:mcmc}

A Bayesian MCMC scan (32 walkers, 5000 steps, emcee, fast JIT solver) over
$(m_\chi, m_\phi, \alpha)$ against 9 astrophysical observables identifies
the viable island:
\begin{align}
  m_\chi &\in [10,\,100]\,{\rm GeV} \\
  m_\phi &\in [7.6,\,14.8]\,{\rm MeV} \\
  \alpha  &\in [5\times10^{-4},\,5\times10^{-3}]
\end{align}

122 benchmark points pass all cuts simultaneously.

\section{Astrophysical Confrontation}
\label{ch1:astro}

[Section to be expanded. Key results summarized in Table~\ref{tab:ch1_pheno}.]

\begin{table}[h!]
\centering
\caption{Phenomenological summary for named benchmark points.}
\label{tab:ch1_pheno}
\begin{tabular}{lcccc}
\toprule
Observable & BP1 & BP9 & MAP & Constraint \\
\midrule
Fornax GC survival & 12/15 & 13/15 & \textbf{14/15} & all 5 survive \\
$r_{\rm core}$ Fornax [pc] & 449 & 541 & \textbf{829} & $\gtrsim 500$ \\
Fornax $\chi^2$ (Jeans+OM) & 35 & — & \textbf{2.3} & — \\
Read+2019 bound & \checkmark & \textsf{FAIL} & $\checkmark$ (margin +0.05) & $\sigma/m < 1.5$ cm$^2$/g \\
RAR scatter & \textbf{0.122} dex & — & 0.28 dex & 0.177 dex (obs.) \\
Sommerfeld $S_0$ at $T_{fo}$ & 1.003 & 1.010 & 1.025 & $< 2.5\%$ \\
\bottomrule
\end{tabular}
\end{table}


%%=================================================================
%%=================================================================
\part{Path-Integral Foundations: From $\mathcal{L}$ to $H_0$}
\label{part:pathintegral}
%%=================================================================
%%=================================================================

%%=================================================================
%% CHAPTER 2: Path-Integral Foundations — From L to H0
%% Source: The_Lagernizant_integral_SIDM/ (within Secluded-Majorana-SIDM)
%%=================================================================

\section{Overview: Seven Layers}
\label{ch2:overview}

This chapter derives all results of Chapter~1 from the generating functional
$Z[J,\eta,\bar\eta]$ of the path integral. Seven layers connect the Lagrangian
to physical observables:

\begin{table}[h!]
\centering
\caption{The seven path-integral layers.}
\begin{tabular}{clll}
\toprule
Layer & Computation & Method & Output \\
\midrule
1 & Free propagators $\Delta_F$, $S_F$ & Gaussian integration & Feynman rules \\
2 & Yukawa potential $V(r)$ & Tree-level (t+u) & $V = -\alpha e^{-m_\phi r}/r$ \\
3 & Cross section $\sigmaT(v)$ & Fluctuation determinant & VPM $\equiv$ exact \\
4 & CW potential $V_{CW}(\theta)$ & One-loop Coleman-Weinberg & $A_4$ vacuum \\
5 & Thermal potential $V_{eff}(\theta,T)$ & Matsubara / finite-$T$ & Freeze-out stability \\
6 & Euclidean action $S_E$ & Instanton & $\theta_{A_4}$ stable ($S_E \sim 10^{121}$) \\
7 & Relic density $\Omega h^2$ & $\langle\sigma v\rangle$ + Boltzmann & $\Omega h^2 = 0.120$ \\
DE & $H_0$ from $V_{eff}(\sigma_0)$ & Dark QCD misalignment & $H_0 = 67.4$ km/s/Mpc \\
\bottomrule
\end{tabular}
\end{table}

\section{Layer 1: Free Propagators}
\label{ch2:propagators}

The Lagrangian:
\begin{equation}
  \Lc = \frac{1}{2}\bar{\chi}(i\slashed{\partial} - \mchi)\chi
      + \frac{1}{2}(\partial\phi)^2 - \frac{1}{2}\mphi^2\phi^2
      - \frac{y}{2}\bar{\chi}\chi\phi
\end{equation}

Gaussian integration over free fields yields:
\begin{align}
  \Delta_F(q^2) &= \frac{i}{q^2 - \mphi^2 + i\varepsilon} \\
  S_F(p)        &= \frac{i(\slashed{p} + \mchi)}{p^2 - \mchi^2 + i\varepsilon}
\end{align}

Physical scales: range $r_0 = 1/\mphi$, de Broglie length
$\lambda_{dB} = 4/(m_\chi v)$. Regime: resonant when $\lambda_{dB}/r_0 < 2$.

\section{Layer 2: Tree-Level Yukawa Potential}
\label{ch2:tree}

Two vertices $\times$ one propagator at momentum transfer $\mathbf{q}$:
\begin{equation}
  i\mathcal{M} = \left(-\frac{iy}{2}\right)^2 \frac{i}{-|\mathbf{q}|^2 - \mphi^2}
  \quad\xrightarrow{\rm NR,\;Fourier}\quad
  V(r) = -\frac{\alpha}{r}e^{-\mphi r}
\end{equation}

Born cross section (valid only for $\lambda \ll 1$):
\begin{equation}
  \sigma_T^{\rm Born} = \frac{16\pi\alpha^2}{m_\chi^2 v^4}\,f(\beta),
  \qquad f(\beta) = \ln(1+\beta^2) - \frac{\beta^2}{1+\beta^2}
\end{equation}

At benchmark couplings, Born overestimates VPM by factors 88–135$\times$.
Full VPM is essential.

\section{Layer 3: VPM as Fluctuation Determinant (Exact)}
\label{ch2:vpm_pi}

\textbf{Central theorem of this work.} The scattering path integral evaluated
around the classical two-body density $n_\chi(r)$ gives:
\begin{equation}
  Z_{\rm scatter} = e^{iS_{cl}} \cdot
  \left[\det\!\left(-\partial^2 - \mphi^2 - y^2 n_\chi(r)\right)\right]^{-1/2}
\end{equation}

The fluctuation determinant evaluates exactly via:
\begin{equation}
  \ln\det = {\rm Tr}\ln = \sum_\ell (2\ell+1)\ln e^{2i\delta_\ell}
           = 2i\sum_\ell(2\ell+1)\delta_\ell
\end{equation}

Hence the cross section:
\begin{equation}
  \sigmaT = \frac{2\pi}{k^2}\sum_\ell w_\ell(2\ell+1)\sin^2\delta_\ell
\end{equation}

\textbf{The VPM phase shifts $\delta_\ell$ \emph{are} the path integral.}
No approximation is made. The ODE
$d\delta_\ell/dx = \ldots$ is the exact quantum mechanical one-loop result.

\section{Layer 4: Coleman-Weinberg Potential}
\label{ch2:cw}

Integrating out the Majorana fermion (2 degrees of freedom) at one loop:
\begin{equation}
  V_{CW}(\theta) = -\frac{2}{64\pi^2}M_{eff}^4(\theta)
  \left[\ln\frac{M_{eff}^2(\theta)}{\mu^2} - \frac{3}{2}\right]
\end{equation}

with $M_{eff}(\theta) = m_\chi\sqrt{\cos^2\theta + \sin^2\theta/9}$ (with $A_4$ symmetry).

Result: $V_{CW}$ minimizes at $\theta = \pi/2$ (pseudoscalar preferred).
The $A_4$ symmetry fixes $\theta = \arcsin(1/3)$ as a discrete vacuum.

\section{Layer 5: Finite-$T$ Matsubara Potential}
\label{ch2:matsubara}

\begin{equation}
  V_{eff}(\theta, T) = V_{CW}(\theta) + \frac{T^4}{2\pi^2}\int_0^\infty dp\,p^2
  \left[\ln(1 - e^{-E_\phi/T}) - 2\ln(1 + e^{-E_\chi/T})\right]
\end{equation}

At $T \gg m_\chi$: thermal corrections are dominant and push $\theta \to 0$.
At $T = T_{fo}$: $V_T/V_{CW} \sim 10^{-4}$ — CW minimum governs freeze-out.
No thermal attractor at $\theta \sim 2$ rad (PI-1, confirmed).

\section{Layer 6: Instanton Stability}
\label{ch2:instanton}

For $m_\sigma \sim H_0$ and $f \sim 0.24\,M_{\rm Pl}$:
\begin{equation}
  S_E \sim \left(\frac{f}{m_\sigma}\right)^2 \sim 10^{121}
  \qquad\Rightarrow\qquad
  \Gamma_{\rm tunnel} = e^{-S_E} \approx 0
\end{equation}

The $A_4$ vacuum is stable to tunneling for $\sim 10^{121}$ Hubble times.

\section{Layer 7: Relic Density from the Path Integral}
\label{ch2:relic}

The $s$-wave annihilation amplitude evaluated from the path integral:
\begin{equation}
  \langle\sigma v\rangle_s = \frac{\pi\alpha^2}{4m_\chi^2}
\end{equation}

Solving the Boltzmann equation $dY/dx = \ldots$ gives $\Omega h^2 = 0.120$
for each $(m_\chi, \alpha_{\rm relic})$ in the viable island. The same $\alpha$
that determines SIDM uniquely sets the relic abundance.

\section{The Hubble Constant from Dark QCD}
\label{ch2:h0}

The full derivation chain:
\begin{equation}
  \mathcal{L}_{SIDM} \xrightarrow{1\text{-loop}} V_{CW}(\theta)
  \xrightarrow{A_4} \theta_{relic}
  \xrightarrow{\rm dark\,QCD} \Omega_\sigma = 0.69
  \xrightarrow{\rm Friedmann} H_0 = 67.4\;\text{km/s/Mpc}
\end{equation}

Dark QCD confinement scale: $\Lambda_d \sim 2.05\,\text{meV}$.
Dark axion mass: $m_\sigma = \Lambda_d^2/f \sim H_0$ (GMOR relation).
Dark energy density from misalignment:
\begin{equation}
  \Omega_\sigma = \frac{f^2\theta_i^2 H_0^2}{2\rho_{crit}} \approx 0.69
  \quad{\rm for}\quad \theta_i \sim 2\,{\rm rad},\;\; f \sim 0.24\,M_{\rm Pl}
\end{equation}

$H_0 = 67.4$~km/s/Mpc is obtained from $V_{eff}(\sigma_0) = \rho_\Lambda$
via the Friedmann equation. This prediction is \emph{parameter-free} once
the SIDM+relic constraints fix $\theta_{relic} = 19.47^\circ$.


%%=================================================================
%%=================================================================
\part{Dark Electromagnetic Duality: Dark Energy as T-Breaking}
\label{part:duality}
%%=================================================================
%%=================================================================

%%=================================================================
%% CHAPTER 3: Dark Electromagnetic Duality
%% Dark Energy as the T-Breaking Component of Dark Matter
%% Source: dark-energy-T-breaking/ (within Secluded-Majorana-SIDM)
%%=================================================================

\section{The Central Analogy}
\label{ch3:analogy}

In classical electromagnetism, a static charge produces only an electric field
$\vec{E}$. Motion reveals a second component — the magnetic field $\vec{B}$ —
which is not a new entity but a different projection of the same tensor $F_{\mu\nu}$.
The critical property: $\vec{B}$ is T-odd while $\vec{E}$ is T-even.

We propose the dark sector obeys the same structure:

\begin{table}[h!]
\centering
\caption{Dark Electromagnetic Duality correspondence.}
\begin{tabular}{ll}
\toprule
Electromagnetism & Dark Sector \\
\midrule
$\vec{E}$ (T-even) & $y_s = y\cos(\sigma/f)$ — attraction, SIDM clustering \\
$\vec{B}$ (T-odd) & $y_p = y\sin(\sigma/f)$ — repulsion, DE acceleration \\
$F_{\mu\nu}$ & $\mathcal{Y} = y_s + iy_p = ye^{i\gamma^5\sigma/f}$ \\
Lorentz boost mixes $E \leftrightarrow B$ & $\sigma$ field rotates $y_s \leftrightarrow y_p$ \\
$\vec{B}$ does no work ($\vec{F}\perp\vec{v}$) & $\gamma^5$ prevents $y_p$ from contributing to mass \\
Like charges repel electrically & $y_s$ causes DM–DM attraction (SIDM) \\
Parallel currents attract magnetically & $y_p$ causes vacuum repulsion (DE) \\
\bottomrule
\end{tabular}
\end{table}

\textbf{Thesis}: Dark energy is the T-violating component of the dark matter
interaction, made manifest through CP violation in the Majorana sector —
exactly as magnetism is the T-violating component of electrostatics,
made manifest through relative motion.

\section{The Extended Lagrangian}
\label{ch3:lagrangian}

Starting from the SIDM Lagrangian (Chapter~1), we promote the CP phase to a
dynamical field $\sigma$ — a dark pseudo-Goldstone boson:
\begin{equation}
  y_s + iy_p \;\longrightarrow\; ye^{i\gamma^5\sigma/f}
  = y\cos\!\left(\frac{\sigma}{f}\right) + iy\sin\!\left(\frac{\sigma}{f}\right)\gamma^5
\end{equation}

The full unified Lagrangian:
\begin{equation}
  \boxed{\Lc = \underbrace{\frac{1}{2}\bar{\chi}(i\slashed{\partial} - \mchi)\chi
    + \frac{1}{2}(\partial\phi)^2 - V_0(\phi)}_{\rm SIDM\;(Ch.1)}
    - \frac{1}{2}\bar{\chi}\!\left[y\cos\!\tfrac{\sigma}{f}
      + iy\sin\!\tfrac{\sigma}{f}\gamma^5\right]\!\chi\,\phi
    + \underbrace{\frac{1}{2}(\partial\sigma)^2 - V_{\rm bare}(\sigma)}_{\rm dark\;axion}}
\end{equation}

Parameters: $y$ fixed by SIDM+relic (Chapter~1); $f \sim 0.2\,M_{\rm Pl}$ from
$V_\sigma \sim \rho_\Lambda$ (see §\ref{ch3:scale}); all other parameters determined.

\section{The Universal Relic Angle}
\label{ch3:thetarel}

Two simultaneous constraints:
\begin{align}
  \text{SIDM (Ch.1):} \quad & \alphass = \alpha \\
  \text{Relic (Ch.1):} \quad & \alphass\alphapp = \frac{\alpha^2}{8}
\end{align}

These fix $\alphapp = \alpha/8$, hence:
\begin{equation}
  \boxed{\thetarel = \arctan\!\sqrt{\frac{\alphapp}{\alphass}}
         = \arctan\!\frac{1}{\sqrt{8}} = 19.47^\circ}
\end{equation}

This angle is \emph{universal} — independent of $(m_\chi, m_\phi, \alpha)$.
It is the unique equilibrium where the same coupling
simultaneously satisfies DM self-interactions and the observed relic abundance.

\section{Scale of the Dark Energy Potential}
\label{ch3:scale}

The effective potential at $\theta = \thetarel$:
\begin{align}
  V_\sigma &\approx V_{CW}(\thetarel) \cdot \left(\frac{f}{M_{\rm Pl}}\right)^4 \\
  V_\sigma &\approx \rho_\Lambda \quad\Leftrightarrow\quad f \approx 0.2\,M_{\rm Pl}
\end{align}

Numerically, for the MAP benchmark:
\begin{equation}
  f_{\rm MAP} = \frac{M_{\rm Pl}}{(V_{CW}/\rho_\Lambda)^{1/4}}
              \approx 6.6\times10^{17}\,{\rm GeV} = 0.27\,M_{\rm Pl}
\end{equation}

The sub-Planckian value $f \sim 0.2$–$0.3\,M_{\rm Pl}$ is natural in string
compactifications and aligns with the weak gravity conjecture.

\section{Fifth Force and Screening}
\label{ch3:fifthforce}

The $\sigma$ field mediates a fifth force with coupling:
\begin{equation}
  \beta = \frac{M_{\rm Pl}}{f} \approx 5.0 \quad \text{(universal for all BPs)}
\end{equation}

Existing Planck+BAO constraints give $\beta < 0.066$ for an unscreened
coupled quintessence field. Resolution requires:
(a) Chameleon screening in dense environments, or
(b) $\sigma$ coupling only to dark matter (secluded),
    not to baryons — consistent with the secluded dark sector hypothesis.

Since $\phi$ has no tree-level SM coupling (§\ref{ch1:model}), $\sigma$ — being
a phase of the $\chi$–$\phi$ system — naturally inherits this seclusion.
Baryonic fifth-force constraints do not directly apply.

\section{Observational Signatures}
\label{ch3:signatures}

\begin{enumerate}
  \item \textbf{Equation of state}: The dark axion potential predicts
        $w_{\rm DE} = -1 + \mathcal{O}(\theta_i^2/3)$ near $\theta_i \sim 2$ rad.
        Current DESI constraint ($w_0 = -0.827 \pm 0.063$) is consistent at
        $\sim 1.6\sigma$ from the model prediction ($w_0 \approx -0.727$).

  \item \textbf{Hubble tension}: $H_0 = 67.4$ km/s/Mpc (from Chapter~2)
        is consistent with CMB measurements and does not worsen the H$_0$ tension
        relative to $\Lambda$CDM.

  \item \textbf{CP violation in SIDM}: The $y_s/y_p$ ratio is in principle
        measurable through indirect detection asymmetries at
        $\mathcal{O}(\sin^2\theta_{relic}) \approx 10\%$ level.

  \item \textbf{Null prediction}: No baryonic fifth force.
        No cross-coupling between $\sigma$ and SM particles at tree level.
\end{enumerate}

\section{Summary: From SIDM to Dark Energy in One Lagrangian}
\label{ch3:summary}

The three parts of this paper are not separate models. They are three
projections of one Lagrangian:

\begin{itemize}
  \item \textbf{T-even projection} ($y_s$, $\phi$): DM self-interactions,
        velocity-dependent cross sections, Fornax cores, relic density.
        This is Chapter~1.
  \item \textbf{Path-integral structure}: VPM = fluctuation determinant,
        CW vacuum, Matsubara freeze-out, dark QCD $\to H_0$.
        This is Chapter~2.
  \item \textbf{T-odd projection} ($y_p$, $\sigma$): vacuum energy, DE equation
        of state, universal relic angle. This is Chapter~3.
\end{itemize}

The unification is not formal — the same coupling $y$, the same mass $m_\chi$,
and the same mediator mass $m_\phi$ that explain the rotation curves of SPARC
galaxies and the survival of Fornax globular clusters also set the scale of
dark energy and predict $H_0 = 67.4$ km/s/Mpc.


%%=================================================================
\appendix
%%=================================================================

\section{VPM as Fluctuation Determinant: Proof}
\label{app:vpm_proof}

The key identity proved in Chapter~2 (PI-6):
\begin{equation}
  \ln \det\!\left[-\partial^2 - m_\phi^2 - y^2 n_\chi(r)\right]^{-1/2}
  = i\sum_\ell (2\ell+1)\delta_\ell
\end{equation}
where $\delta_\ell$ are the VPM phase shifts. This establishes that the
Variable Phase Method is not an approximation but the \emph{exact} evaluation
of the one-loop path integral over the mediator field $\phi$ in the presence
of the two-body DM density $n_\chi(r)$.

\section{Universal Relic Angle: Derivation}
\label{app:relic_angle}

From the two constraints:
\begin{align}
  \text{SIDM:} & \quad \alphass = \alpha \\
  \text{Relic:} & \quad \alphass \alphapp = \alpha^2/8
\end{align}
it follows immediately that $\alphapp = \alpha/8$, hence:
\begin{equation}
  \thetarel = \arctan\!\sqrt{\alphapp/\alphass} = \arctan(1/\sqrt{8}) = 19.47^\circ
\end{equation}
This angle is \emph{universal} — independent of $(m_\chi, m_\phi, \alpha)$.

\section{Instanton Suppression of $\theta_{A_4}$}
\label{app:instanton}

For a dark axion with $m_\sigma \sim H_0$ and $f \sim 0.24\,M_{\rm Pl}$:
\begin{equation}
  S_E \sim \left(\frac{f}{m_\sigma}\right)^2 \sim 10^{121}
  \qquad \Rightarrow \qquad
  \Gamma_{\rm tunnel} \sim e^{-S_E} \approx 0
\end{equation}
The $A_4$ vacuum is stable to tunneling on timescales
$10^{121} \times t_{\rm Hubble}$ — effectively eternal.

\section{Benchmark Points}
\label{app:benchmarks}

\begin{table}[h!]
\centering
\caption{Named benchmark points from the MCMC posterior (Chapter~1).}
\begin{tabular}{lcccccc}
\toprule
BP & $m_\chi$ [GeV] & $m_\phi$ [MeV] & $\alpha$ & $\lambda$ & $\Omega h^2$ & Source \\
\midrule
BP1  & 20.69 & 11.34 & $1.048\times10^{-3}$ & 1.91 & 0.120 & relic \\
BP9  & 37.93 & 12.98 & $1.858\times10^{-3}$ & 5.43 & 0.120 & relic \\
MAP  & 94.07 & 11.10 & $5.734\times10^{-3}$ & 48.6 & 0.076 & MCMC \\
\bottomrule
\end{tabular}
\end{table}

%%=================================================================
\bibliographystyle{unsrt}
\bibliography{references}

\end{document}
